\documentclass[titlepage,fontsize=10.5 pt]{jlreq}
% beamerならばaspectratio=169をoptionに追加し、jlreq -> beamer

\usepackage{amssymb,amsmath,amsthm,amscd}
\usepackage{ascmac}
\usepackage{bm}
\usepackage{color}
\usepackage{derivative}
\usepackage{enumitem}%beamer metropolisを使う場合コメントアウト
\usepackage{fancybox}
\usepackage{frame,framed}
\usepackage{float}
\usepackage{graphicx}
\usepackage{mathtools}
\usepackage{physics2}
\usepackage{siunitx}
\usepackage{tabularray}
\usepackage{tcolorbox}
\usepackage{tikz}
\usepackage{url}
\usepackage{titlepic} %beamer metropolisを使う場合コメントアウト
\usepackage{wrapfig}

\usepackage[hang,small,bf]{caption}
\usepackage[subrefformat=parens]{subcaption}
\captionsetup{compatibility=false}

\renewcommand{\theequation}{\thesection.\arabic{equation}}
\makeatletter
\@addtoreset{equation}{section}
\makeatother

\usephysicsmodule{ab,braket,doubleprod,xmat}
\DefTblrTemplate{contfoot-text}{normal}{（次ページに続く）}
\SetTblrTemplate{contfoot-text}{normal}
\DefTblrTemplate{conthead-text}{normal}{（続き）}
\SetTblrTemplate{conthead-text}{normal}

%---for beamer metropolis---%
%---If you use beamer metropolis,remove % below sentense---%
% \usepackage[haranoaji,deluxe]{luatexja-preset}% フォント指定
% \usepackage{luatexja}
% \renewcommand{\kanjifamilydefault}{\gtdefault}% 既定をゴシック体に
% \usetheme[block=fill]{metropolis}
%---for beamer metropolis end---%

\begin{document}

\title{電磁気学}
\author{平山大知/HiBand Beige}
\date{\today}
\maketitle
\tableofcontents
\clearpage
\section{はじめに}
以下の文書は、筆者自身が電磁気学の演習をしながらまとめることを意図して電磁気学についてをまとめたものである。構成としてはマクスウェル方程式とローレンツの式をまず導入し、以降において種々の電磁気に係る現象を説明することをこころみる。また回路理論についても電磁気学との対応を説明しながら記述し、等価回路によるモデル化などについても説明する。
\section{ベクトル解析とマクスウェル方程式}
\subsection{ナブラ}
第一にナブラ$\nabla$と呼ばれるものを導入する。特に3次元のものについて以下のようになる
\begin{align}
    \nabla &= \begin{pmatrix}
        \dfrac{\partial}{\partial x} , \dfrac{\partial}{\partial y} , \dfrac{\partial }{\partial z} \\
    \end{pmatrix}\label{eq.nabla}
\end{align}
$n$次元に一般化した場合は次のように書くことができる
\begin{align}
    \nabla &= \sum_{i=1}^{n} e_{i} \pdv{}{x_{i}}
\end{align}
なおこのとき$e_{i}$は空間の各基底の単位ベクトルであり、$x_{i}$は空間の各基底である。$n=3$で先の3次元空間、$n=2$で2次元平面、$n=1$では一次元直線上のことを示すことになる。また1次元のみのナブラは$\dfrac{\partial}{\partial x}$であり、まさしく位置についての微分をするものである。

これはみての通りベクトル量であり、ベクトルとして扱う。そしてこれは対象の空間についての微分を行うものになる。
\subsection{場}
場という概念について説明する必要がある、まずそのためには近接作用論と遠隔作用論の説明をする必要がある。

近接作用論とは、物体が受ける作用や力は物体の直近にあるなにかしらによるものであると考える方法である。対して遠隔作用論とは、物体が受ける作用や力は物体がどれだけ隔てられている状況でも直接的になされるものであると考える方法である。

ニュートンによる力学などでは遠隔作用論が支持されていたが、その後に電磁気学が発展するにつれて遠隔作用論では説明ができない事象が発生したため、その説明が可能である近接作用論が支持されるようになった。

近接作用論では、例えば二物体間の相互作用ならば
\end{document}